\subsection{LoRA(Low-Rank Adaptation):低ランク適応の基礎}

問題設定と動機について述べる．
事前学習済みの大規模モデルを下流タスクへ適応させる際,全パラメータの微調整(full fine-tuning)は計算量と記憶容量の双方で高コストである.加えて継続学習の文脈では,本体重みを直接更新すると過去に獲得した知識を上書きし,破滅的忘却を誘発する.LoRA \cite{lora} は,事前学習重み $W_0 \in \mathbb{R}^{d \times k}$ を凍結したまま,その更新量 $\Delta W$ を低ランク行列の積として表現することで,適応に要する学習可能パラメータを大幅に削減する手法である.本研究のように本体(MM-Grounding DINO)を凍結しドメイン毎の差分のみを学習する枠組みにおいて,LoRA はその最も基礎的な構成要素となる.

中核アイデアについて述べる．
LoRA の出発点は,事前学習モデルを下流タスクへ適応させる際の重み更新 $\Delta W$ が低い内在ランク(intrinsic rank)を持つという経験的観察 \cite{intrinsicdim} である.すなわち $\Delta W$ をフルランクで学習する必要はなく,ランク $r \ll \min(d,k)$ の低ランク近似で十分な適応が可能である,という仮説に立つ.

定式化について述べる．
入力 $x$ に対する線形層の順伝播を考える.LoRA は更新量 $\Delta W$ を二つの行列 $B \in \mathbb{R}^{d \times r}$ と $A \in \mathbb{R}^{r \times k}$ の積で表し,
\begin{equation}
W = W_0 + \Delta W = W_0 + BA, \quad \mathrm{rank}(BA) \leq r,
\end{equation}
と再パラメータ化する.順伝播は
\begin{equation}
h = W_0 x + \Delta W x = W_0 x + B A x
\end{equation}
となる.学習時には $W_0$ を凍結し,$A$ と $B$ のみを最適化する.初期化は $A \sim \mathcal{N}(0,\sigma^2)$,$B = \mathbf{0}$ とし,学習開始時点で $\Delta W = BA = \mathbf{0}$ となるよう設計する.これにより適応の初期状態が事前学習モデルと一致する.推論時にはスケーリング係数 $\alpha$ を用いて
\begin{equation}
h = W_0 x + \frac{\alpha}{r} B A x
\end{equation}
と適用する.$\alpha / r$ は実効的な学習率と更新の大きさを調整する役割を担い,$r$ を変更しても更新量のスケールが概ね保たれるよう設計されている.

アルゴリズム手順について述べる．
\begin{enumerate}
\item 適応対象の線形層(注意機構の $W_q, W_v$ など)を選択する.
\item 各対象層に対し $B = \mathbf{0}$,$A$ をランダム初期化し,$W_0$ を凍結する.
\item 下流タスクの損失 $\mathcal{L}$ を $A, B$ についてのみ逆伝播し更新する.
\item 推論時は $W_0 + (\alpha/r) BA$ を統合して用いる.統合後は追加の推論コストが生じない.
\end{enumerate}

主要な実験設定と数値について述べる．
原論文 \cite{lora} では GPT-3 175B などに対し,学習可能パラメータを全微調整の約 $1/10000$ 程度まで削減しつつ,GLUE 等のベンチマークで全微調整に匹敵または上回る性能を報告している(具体的な削減率はモデル・タスクに依存,詳細値は要確認).推論レイテンシの増加が無い点が,アダプタ系手法に対する優位点として強調されている.

長所短所について述べる．
長所は,(i) 学習・保存コストが低い,(ii) 統合により推論オーバヘッドが無い,(iii) 本体凍結により事前学習知識を直接破壊しない,点である.短所は,(i) ランク $r$ を固定的に与える必要があり層毎の最適配分ができない,(ii) 単体では複数タスクを逐次学習する際のタスク間干渉を抑制する機構を持たない,点である.後者は継続学習において致命的であり,本章の後続手法はいずれもこの干渉抑制を主眼とする.

本研究への含意について述べる．
本研究では本体を凍結しドメイン毎に低ランク差分を学習する点で LoRA を基礎とするが,単純な LoRA の逐次適用では新ドメイン学習時に過去ドメインおよび COCO(task0)の知識が劣化する.したがって LoRA の低ランク構造を保ちつつ,更新部分空間を過去勾配空間と直交化する後続の枠組み(特に InfLoRA \cite{inflora},O-LoRA \cite{olora})が,提案手法の直接の出発点となる.

\subsection{AdaLoRA:特異値による動的ランク配分}

問題設定と動機について述べる．
LoRA はすべての対象層に同一のランク $r$ を割り当てるが,実際には層やモジュールによって適応に要する自由度は大きく異なる.重要な層に過小なランクを与えると適応不足を招き,重要でない層に過大なランクを与えるとパラメータの浪費と過学習を招く.AdaLoRA \cite{adalora} は,固定予算(総ランク数)を層間で重要度に応じて動的に再配分することでこの問題に対処する.

中核アイデアについて述べる．
更新量 $\Delta W$ を特異値分解(SVD)の形に明示的に再パラメータ化し,各特異三つ組(左特異ベクトル・特異値・右特異ベクトル)を重要度で評価する.重要度の低い特異値を学習過程で漸進的に零化(prune)することで,各層の実効ランクを自動的に決定する.

定式化について述べる．
AdaLoRA は更新量を
\begin{equation}
\Delta W = P \Lambda Q,
\end{equation}
と分解する.ここで $P \in \mathbb{R}^{d \times r}$ は左特異ベクトル,$Q \in \mathbb{R}^{r \times k}$ は右特異ベクトルに対応し,$\Lambda = \mathrm{diag}(\lambda_1, \dots, \lambda_r)$ は特異値の対角行列である.$P, Q$ には直交性を促す正則化
\begin{equation}
\mathcal{R}(P,Q) = \| P^\top P - I \|_F^2 + \| Q Q^\top - I \|_F^2
\end{equation}
を課し,厳密な SVD を毎反復計算する代償なしに近似的な特異構造を維持する.第 $i$ 特異三つ組の重要度スコア $s_i$ は,特異値 $\lambda_i$ とそれに付随する特異ベクトルの感度(勾配と値の積の移動平均)に基づいて算出され,大域的予算スケジューラのもとで重要度下位の三つ組の $\lambda_i$ を $0$ に設定する.総損失は
\begin{equation}
\mathcal{L}_{\text{total}} = \mathcal{L}_{\text{task}} + \gamma \, \mathcal{R}(P,Q)
\end{equation}
の形をとる($\gamma$ は正則化係数).

アルゴリズム手順について述べる．
\begin{enumerate}
\item 各層の更新を $P \Lambda Q$ として初期化する.
\item 順伝播・逆伝播でタスク損失と直交正則化を最小化する.
\item 各特異三つ組の感度に基づく重要度 $s_i$ を移動平均で更新する.
\item 予算スケジューラに従い,重要度下位の特異値を零化して実効ランクを縮小する.
\item 学習が進むにつれ目標予算へ収束させる.
\end{enumerate}

主要な実験設定と数値について述べる．
原論文 \cite{adalora} は GLUE,SQuAD,XSum,CNN/DailyMail で評価し,特に低予算条件下で LoRA を一貫して上回ることを報告している(具体的なスコアはタスク依存,詳細値は要確認).

長所短所について述べる．
長所は重要度に応じたランクの自動配分により限られた予算を効率的に活用できる点である.短所は,(i) 重要度推定と零化のための追加計算・ハイパパラメータが必要,(ii) 本来は単一タスクの効率化手法であり,継続学習におけるタスク間干渉抑制を直接の目的としない点である.

本研究への含意について述べる．
AdaLoRA の特異値に基づく重要度評価という考え方は,本研究において各ドメインの更新部分空間の有効次元を見積もる際や,DualGPM \cite{inflora} で部分空間を SVD により拡張・縮約する際の指針として援用しうる.ただし提案手法の主目的は干渉抑制であるため,AdaLoRA は補助的な位置づけにとどまる.

\subsection{プロンプト系継続学習:L2P,DualPrompt,CODA-Prompt}

問題設定と動機について述べる．
プロンプト系継続学習は,事前学習済み Vision Transformer を凍結し,学習可能な少数の「プロンプト」ベクトルを入力に付与することで逐次的なタスク適応を実現する枠組みである.本体を更新しないため事前学習知識の保持に優れ,リハーサル(過去データの保存)を必要としない点で,本研究の本体凍結方針と思想を共有する.

中核アイデアについて述べる．
L2P \cite{ltp} はプロンプトの「プール」を導入し,入力に応じてキー照合により関連するプロンプトを選択する.DualPrompt \cite{dualprompt} はタスク共通の知識を担う general prompt と,タスク固有の知識を担う expert prompt を分離し,異なる層へ付与する.CODA-Prompt \cite{codaprompt} はこれらの離散的なプロンプト選択を,入力条件付きの連続的な重み付き合成へ置き換える.

定式化について述べる．
CODA-Prompt \cite{codaprompt} では,プロンプト成分の集合 $\{P_m\}_{m=1}^{M}$,対応するキー $\{k_m\}$ およびアテンション $\{a_m\}$ を用意する.入力 $x$ の特徴量 $q(x)$ に対し,各成分の重みを
\begin{equation}
\alpha_m = \cos\!\big(q(x) \odot a_m,\; k_m\big)
\end{equation}
で計算し($\odot$ は要素積,$\cos$ は余弦類似度),入力条件付きのプロンプトを
\begin{equation}
P = \sum_{m=1}^{M} \alpha_m \, P_m
\end{equation}
として合成する.干渉を抑えるため,プロンプト成分とキー,アテンションには直交性制約
\begin{equation}
\mathcal{L}_{\text{ortho}} = \| B^\top B - I \|_F^2
\end{equation}
を課す($B$ は $P_m$,$k_m$,$a_m$ を列に並べた行列).新タスク学習時は過去タスクに割り当てた成分を凍結し,新規成分のみを学習することで,選択の鋭さを保ちつつ過去成分への干渉を避ける.

アルゴリズム手順について述べる．
\begin{enumerate}
\item 凍結 ViT の上にプロンプト成分・キー・アテンションを定義する.
\item 入力特徴 $q(x)$ とキーの照合で各成分の重み $\alpha_m$ を計算する.
\item 重み付き和で入力条件付きプロンプト $P$ を合成し,中間層に挿入する.
\item タスク損失に直交正則化 $\mathcal{L}_{\text{ortho}}$ を加えて学習する.
\item 新タスクでは過去成分を凍結し,新規成分のみ学習する.
\end{enumerate}

主要な実験設定と数値について述べる．
これらは主に ImageNet-R,CIFAR-100,DomainNet 等のクラス増分学習で評価される.CODA-Prompt \cite{codaprompt} は L2P・DualPrompt を上回る平均精度を報告している(具体値は設定依存,詳細は要確認).

長所短所について述べる．
長所は本体完全凍結による高い知識保持とリハーサル不要性である.短所は,(i) タスク数増加に伴いプロンプトプールが拡大しうる,(ii) 入力依存の成分選択がタスク識別に近い役割を担い,選択誤りが性能に直結する,(iii) プロンプトは入力空間の摂動であり,検出ヘッドの言語・視覚整合のような構造的更新には適応力が限定的,点である.

本研究への含意について述べる．
CODA-Prompt の直交性に基づく成分合成・過去成分凍結という設計思想は,本研究で採用する更新部分空間の直交化と整合する.ただし本研究は物体検出という密な構造化予測を対象とし,プロンプト摂動よりも重み空間への低ランク更新(LoRA 系)の方が表現力・統合容易性の面で適する.したがってプロンプト系は思想的先行として位置づけ,実装は LoRA 系を基盤とする.

\subsection{InfLoRA:干渉なし低ランク適応}

問題設定と動機について述べる．
継続学習では,新タスクの学習が過去タスクの性能を劣化させる干渉(interference)が中心的課題である.LoRA を逐次適用するだけでは,新タスクの更新 $\Delta W_t = B_t A_t$ が過去タスクの入力部分空間に成分を持ち,過去出力を変化させてしまう.InfLoRA \cite{inflora} は,この干渉を学習開始前に幾何的に排除するよう,更新部分空間そのものを事前設計する手法である.本研究の提案手法(InfLoRA 型 + COCO-task0)が直接基盤とする最重要手法の一つである.

中核アイデアについて述べる．
LoRA の更新 $B_t A_t x$ において,$B_t$ の列空間(更新が及ぶ出力方向)を,(i) 新タスクの勾配(入力)が張る空間 $\mathcal{N}_t$ の内部にあり,かつ (ii) 過去タスクの勾配が張る空間 $\mathcal{M}_{t-1}$ の直交補空間 $\mathcal{M}_{t-1}^{\perp}$ にある,という条件を満たすよう事前に固定する.これにより,$A_t$ をどう学習しても新タスクの更新は過去タスクの入力に直交し,過去出力を一切変化させない(干渉が原理的に零となる).学習対象は $A_t$ のみで,$B_t$ は固定する.

定式化について述べる．
第 $t$ タスクの線形層出力を
\begin{equation}
h = W_0 x + B_t A_t x,
\end{equation}
とする.ここで $B_t \in \mathbb{R}^{d \times r}$ は固定,$A_t \in \mathbb{R}^{r \times k}$ のみ学習する.過去タスクの入力(勾配)部分空間を $\mathcal{M}_{t-1} = \mathrm{Span}(\mathbf{M}_{t-1})$ とし,その直交射影行列を
\begin{equation}
\mathbf{P}_{\mathcal{M}_{t-1}} = \mathbf{M}_{t-1} \mathbf{M}_{t-1}^\top
\end{equation}
と書く.新タスクの入力相関を表す行列を $\mathbf{N}_t$(現タスク特徴の主成分から構成)とするとき,InfLoRA は $B_t$ の列空間を
\begin{equation}
\mathrm{Span}(B_t) \subseteq \mathcal{N}_t \cap \mathcal{M}_{t-1}^{\perp}
\end{equation}
となるよう構成する.具体的には,新タスク特徴を過去部分空間の直交補へ射影し,
\begin{equation}
\widehat{\mathbf{N}}_t = \big(\mathbf{I} - \mathbf{P}_{\mathcal{M}_{t-1}}\big)\,\mathbf{N}_t,
\end{equation}
の主要な特異方向(上位 $r$ 本)を $B_t$ の列とする.このとき任意の $x$ について更新成分は
\begin{equation}
\mathbf{P}_{\mathcal{M}_{t-1}} B_t A_t x = \mathbf{0}
\end{equation}
を満たし,過去タスクの入力方向への影響が消える.

部分空間の維持には DualGPM(Dual Gradient Projection Memory)を用いる.GPM \cite{gpm} は過去入力の相関行列を SVD し,主要方向を基底として蓄積するが,DualGPM はタスク増加に応じて,$\mathcal{M}$(蓄積空間)の拡張と,その直交補 $\mathcal{M}^{\perp}$ の縮約の双方を陽に管理する.第 $t$ タスク後の更新は,新タスク入力 $\mathbf{R}_t$ を既存基底へ射影した残差
\begin{equation}
\mathbf{R}_t^{\text{res}} = \mathbf{R}_t - \mathbf{M}_{t-1}\mathbf{M}_{t-1}^\top \mathbf{R}_t
\end{equation}
を SVD し,エネルギー閾値を超える主要方向を $\mathbf{M}_{t-1}$ に追加して $\mathbf{M}_t$ を得る.これにより部分空間がタスク間で一貫して保たれる.

なお InfLoRA は,各タスクで導入した $B_t A_t$ を学習後に本体へ統合(分岐統合)できるため,タスク数が増えてもパラメータ数および推論コストが一定に保たれる.

アルゴリズム手順について述べる．
\begin{enumerate}
\item 過去部分空間 $\mathbf{M}_{t-1}$ の直交射影 $\mathbf{P}_{\mathcal{M}_{t-1}}$ を用意する.
\item 新タスク特徴の相関 $\mathbf{N}_t$ を直交補へ射影し $\widehat{\mathbf{N}}_t$ を得る.
\item $\widehat{\mathbf{N}}_t$ の上位 $r$ 特異方向で $B_t$ を構成し固定する.
\item $A_t$ のみをタスク損失で学習する($A_t$ は通常零初期化).
\item DualGPM で残差 $\mathbf{R}_t^{\text{res}}$ を SVD し $\mathbf{M}_t$ を更新する.
\item $B_t A_t$ を本体へ統合し,次タスクへ進む.
\end{enumerate}

主要な実験設定と数値について述べる．
InfLoRA \cite{inflora} は ImageNet-R,CIFAR-100,DomainNet を用いたクラス増分学習で評価され,L2P・DualPrompt・CODA-Prompt 等のプロンプト系および LoRA 系手法を上回る平均精度と低い忘却を報告している(具体値は分割・設定に依存するため要確認).CVPR 2024 採択.

長所短所について述べる．
長所は,(i) 干渉を学習前に幾何的に排除するため理論的に明快,(ii) $A_t$ のみ学習で効率的,(iii) 統合によりパラメータ・推論コストが一定,点である.短所は,(i) 部分空間維持(DualGPM)に過去入力統計の蓄積が必要,(ii) $\mathcal{N}_t \cap \mathcal{M}_{t-1}^{\perp}$ の次元がタスク増加とともに枯渇し可塑性が低下しうる,(iii) 完全な直交化は新タスク適応の自由度を制約しうる,点である.

本研究への含意について述べる．
本研究は InfLoRA を直接の骨格とする.具体的には,task0 を COCO とし,COCO の勾配(入力)部分空間 $\mathcal{M}_0$ を初期過去空間として登録することで,以降の各ドメイン適応の更新を $\mathcal{M}_0^{\perp}$ 方向へ制約し,COCO のゼロショット検出性能(ZCOCO)の保持を干渉抑制として実現する.差別化点は,(i) 対象が分類でなく物体検出(MM-Grounding DINO)であり,$\mathbf{N}_t$ や $\mathbf{M}$ の推定を検出器の視覚・言語特徴に対して行う必要がある,(ii) クラス増分でなくドメイン増分かつ task0=COCO という非対称な保持要件を持つ,点にある.部分空間枯渇への対処として,直交制約と緩和(後述の SD-LoRA・O-LoRA 的損失)を組み合わせる余地がある.

\subsection{O-LoRA:直交部分空間学習}

問題設定と動機について述べる．
InfLoRA が更新部分空間を学習前に厳密に直交設計するのに対し,O-LoRA \cite{olora} は過去タスクの勾配部分空間を陽に蓄積・射影することなく,過去 LoRA 行列そのものを過去勾配方向の代理として用い,直交性を損失項として柔らかく強制する.実装が簡潔でリハーサル不要であり,大規模言語モデルの継続学習でゼロショット汎化能力の保持に優れる点で,本研究の ZCOCO 保持要件と思想を共有する.

中核アイデアについて述べる．
LoRA の更新 $\Delta W_t = B_t A_t$ において,$A_t$ の列空間が下流タスクの勾配が張る部分空間を良く近似するという観察に基づく.過去タスクの $A_i$($i < t$)を過去勾配部分空間の代理とみなし,新タスクの $A_t$ をそれらと直交させることで,新旧タスクの更新が異なる部分空間に収まり干渉を抑える.

定式化について述べる．
過去タスク $i$ の LoRA 行列 $A_i$ と新タスク $A_t$ の直交性を測る損失を
\begin{equation}
\mathcal{L}_{\text{orth}} = \sum_{i=1}^{t-1} \big\| A_i^\top A_t \big\|_F^2,
\end{equation}
と定義する($\|\cdot\|_F$ は Frobenius ノルム).$A_i^\top A_t = \mathbf{0}$ のとき両部分空間が直交する.総損失は,新タスクのタスク損失 $\mathcal{L}_{\text{task}}$ と上式に加え,過去 LoRA を固定するための正則化も含めて
\begin{equation}
\mathcal{L}_{\text{total}} = \mathcal{L}_{\text{task}} + \lambda_1 \mathcal{L}_{\text{orth}} + \lambda_2 \mathcal{L}_{\text{l2}},
\end{equation}
の形をとる($\lambda_1, \lambda_2$ は重み係数).学習時は過去タスクの $\{A_i, B_i\}$ を固定し,新タスクの $A_t, B_t$ のみを更新しつつ,$\mathcal{L}_{\text{orth}}$ で過去方向との直交を促す.

アルゴリズム手順について述べる．
\begin{enumerate}
\item 過去タスクの LoRA 行列 $\{A_i\}_{i<t}$ を保持・固定する.
\item 新タスク用 $A_t, B_t$ を初期化する.
\item 各反復で $\mathcal{L}_{\text{task}}$ と直交損失 $\mathcal{L}_{\text{orth}} = \sum_i \|A_i^\top A_t\|_F^2$ を計算する.
\item $\mathcal{L}_{\text{total}}$ を $A_t, B_t$ について逆伝播・更新する.
\item 学習後,$A_t$ を過去集合に追加し次タスクへ進む.
\end{enumerate}

主要な実験設定と数値について述べる．
O-LoRA \cite{olora} は標準的な継続学習ベンチマーク(15タスク・T5 ベース等)で,既存手法を上回る平均精度を報告している.特に注目すべきは,Alpaca で指示チューニングした LLaMA-7B を用いた継続学習後の MMLU ゼロショット評価で,O-LoRA が $33.6\%$,通常の逐次 LoRA(LoRA-CL)が $23.3\%$ であり,直交化により未見タスクへの汎化能力(ゼロショット性能)が顕著に保持されることを実証した点である.EMNLP 2023 Findings 採択.

長所短所について述べる．
長所は,(i) 過去入力統計の蓄積や明示的な部分空間維持を要さず実装が簡潔,(ii) リハーサル不要,(iii) ゼロショット汎化の保持を実証,点である.短所は,(i) $A_t$ を過去勾配の代理とする近似であり厳密な直交は保証されない,(ii) 直交損失はあくまで軟制約のため $\lambda_1$ の調整に敏感,(iii) タスク増加に伴い過去 $A_i$ の保持コストと直交損失の項数が増える,点である.

本研究への含意について述べる．
O-LoRA の「ゼロショット汎化保持」の実証は,本研究の ZCOCO 保持目標に直接対応する強い動機づけを与える.提案手法では InfLoRA の厳密直交設計を骨格としつつ,task0=COCO の $A_0$(あるいは COCO 勾配部分空間)に対する直交損失 $\sum_i \|A_0^\top A_t\|_F^2$ を補助項として加えることで,部分空間枯渇時にも COCO 方向への漏れを軟制約で抑える折衷が考えられる.差別化点は,(i) 言語モデルの MMLU ではなく検出器の COCO ゼロショット mAP を保持対象とする,(ii) 過去タスク間の対称的直交でなく,COCO への非対称な強い保持を課す,点にある.

\subsection{SD-LoRA:大きさと方向の分離}

問題設定と動機について述べる．
多くの LoRA 系・プロンプト系継続学習手法は,タスク数の増加に伴い LoRA/プロンプトのプールを拡張するか,過去データを保持する必要があり,スケーラビリティに課題がある.SD-LoRA \cite{sdlora} は,LoRA 成分を「大きさ(magnitude)」と「方向(direction)」へ分離し,過去に学習した方向を凍結したうえで,大きさと新方向のみを学習することで,プール拡張もリハーサルも不要とする.

中核アイデアについて述べる．
weight-decomposition の発想に基づき,各タスクの LoRA 更新を,正規化された方向成分と,それをスケールするスカラ(大きさ)に分解する.継続学習が進むにつれ過去タスクの方向は固定され,共有の大きさパラメータと新タスク方向のみが更新される.これにより,学習軌跡が全タスクに共通する低損失領域へ収束しやすく,安定性と可塑性の良好な両立が得られると理論・実験の双方で示されている.

定式化について述べる．
第 $t$ タスクまでの累積更新を,各タスクの正規化された方向 $\widehat{B_i A_i} = B_i A_i / \| B_i A_i \|$ と大きさ $s_i \in \mathbb{R}$ を用いて
\begin{equation}
W = W_0 + \sum_{i=1}^{t} s_i \, \frac{B_i A_i}{\| B_i A_i \|}
\end{equation}
と表す.第 $t$ タスク学習時は過去方向 $\{B_i A_i\}_{i<t}$ を凍結し,大きさ $\{s_i\}_{i \leq t}$ と新方向 $B_t A_t$ のみを学習する.これにより新規パラメータの追加はタスクごとに方向一組と少数のスカラに限られ,プールの線形拡張を回避する.推論時は最終的に学習された単一モデルをそのまま用いればよく,タスク ID に基づく成分選択を要しない.

アルゴリズム手順について述べる．
\begin{enumerate}
\item 過去タスクの方向成分 $\{B_i A_i\}_{i<t}$ を凍結する.
\item 新タスク方向 $B_t A_t$ と全タスクの大きさ $\{s_i\}$ を初期化する.
\item タスク損失を $B_t A_t$ と $\{s_i\}$ について最適化する.
\item 学習後,方向 $B_t A_t$ を固定集合へ加える.
\item 推論は最終モデル $W_0 + \sum_i s_i \widehat{B_i A_i}$ を直接用いる.
\end{enumerate}

主要な実験設定と数値について述べる．
SD-LoRA \cite{sdlora} は ImageNet-R,DomainNet,CIFAR-100 等のクラス増分学習で評価され,既存の LoRA・プロンプト系を上回る性能を報告している.加えて,rank reduction を用いる SD-LoRA-RR と知識蒸留を用いる SD-LoRA-KD の二つの省パラメータ変種を提案している(具体値は要確認).ICLR 2025 Oral 採択.

長所短所について述べる．
長所は,(i) プール拡張・リハーサル不要でスケーラブル,(ii) 推論時に成分選択(タスク ID)が不要,(iii) 安定性と可塑性の両立が理論的に裏づけられている,点である.短所は,(i) 過去方向の凍結により方向多様性が時間とともに飽和しうる,(ii) 干渉抑制が直交化のように陽に保証されるわけではなく,主に低損失領域への収束という間接的機構による,点である.

本研究への含意について述べる．
SD-LoRA の大きさ・方向分離とプール非拡張という性質は,ドメインが増加し続ける本研究の設定でパラメータ一定性を保つうえで有用である.提案手法では InfLoRA による直交方向設計と SD-LoRA のスカラ大きさ制御を組み合わせ,COCO(task0)方向の大きさを凍結しつつ新ドメイン方向のみ可塑性を持たせる,という設計が ZCOCO 保持に整合的でありうる.タスク ID 不要の推論は,ドメイン境界が曖昧な検出運用において実務的利点となる.

\subsection{Low-rank Orthogonal Subspaces:Stiefel多様体上の直交部分空間学習}

問題設定と動機について述べる．
本手法 \cite{lowrankorth} は,LoRA 以前の継続学習研究でありながら,InfLoRA や O-LoRA の「タスク毎に互いに直交する低ランク部分空間で学習する」という思想の直接の先行に位置づけられる.各タスクを互いに直交する部分空間で学習すれば,異なるタスクから来る勾配が直交し,一方の学習が他方を破滅的に干渉しないという原理を,Stiefel 多様体上の最適化として定式化した.

中核アイデアについて述べる．
ネットワークの一部の線形写像を等長写像(isometry,直交行列)に制約し,各タスクの表現を互いに直交する低ランク部分空間へ割り当てる.直交行列の集合は Stiefel 多様体をなすため,学習をこの多様体上の制約付き最適化として解く.直交性により,タスク間の勾配干渉が幾何的に抑制される.

定式化について述べる．
タスク $i$ の表現を低ランク部分空間 $\mathcal{U}_i = \mathrm{Span}(\mathbf{U}_i)$,$\mathbf{U}_i \in \mathbb{R}^{d \times r}$ に割り当て,タスク間に直交条件
\begin{equation}
\mathbf{U}_i^\top \mathbf{U}_j = \mathbf{0} \quad (i \neq j), \qquad \mathbf{U}_i^\top \mathbf{U}_i = \mathbf{I}_r,
\end{equation}
を課す.$\mathbf{U}_i^\top \mathbf{U}_i = \mathbf{I}_r$ は $\mathbf{U}_i$ が Stiefel 多様体 $\mathrm{St}(d,r) = \{\mathbf{U} \in \mathbb{R}^{d\times r} : \mathbf{U}^\top \mathbf{U} = \mathbf{I}_r\}$ の元であることを意味する.学習はこの多様体上の Riemannian 最適化で行い,通常勾配 $\nabla \mathcal{L}$ を接空間へ射影した Riemannian 勾配
\begin{equation}
\mathrm{grad}\,\mathcal{L} = \nabla \mathcal{L} - \mathbf{U}\,\mathrm{sym}\!\big(\mathbf{U}^\top \nabla \mathcal{L}\big),
\end{equation}
($\mathrm{sym}(\mathbf{X}) = \tfrac{1}{2}(\mathbf{X} + \mathbf{X}^\top)$)を用い,retraction によって多様体上へ引き戻して更新する.分類には,部分空間構造と相性の良いコサイン類似度に基づく損失を用いる.

アルゴリズム手順について述べる．
\begin{enumerate}
\item 各タスクへ直交する低ランク部分空間 $\mathbf{U}_i$ を割り当てる.
\item 等長(直交)制約を Stiefel 多様体上の制約として設定する.
\item 通常勾配を接空間へ射影し Riemannian 勾配を得る.
\item retraction で多様体上へ更新する.
\item 少量のエピソード記憶(experience replay)と併用して安定化する.
\end{enumerate}

主要な実験設定と数値について述べる．
原論文 \cite{lowrankorth} は,Split CIFAR-100・Split miniImageNet 等で評価し,同量のエピソード記憶を用いる強力な experience replay ベースラインに対して,平均精度で約 $8\%$ の改善,忘却で約 $50\%$ の削減を報告している.NeurIPS 2020 採択.

長所短所について述べる．
長所は,直交部分空間による干渉抑制を多様体最適化として原理的に定式化した点である.短所は,(i) Stiefel 多様体上の最適化は実装・計算が通常の SGD より複雑,(ii) タスク数増加に伴い直交部分空間を確保する余地が枯渇,(iii) 少量のリプレイ記憶に依存する,点である.

本研究への含意について述べる．
本手法は InfLoRA \cite{inflora} の「直交部分空間での更新」の理論的源流であり,本研究の直交化方針の根拠を与える.提案手法は,Stiefel 多様体上の重い最適化を避けつつ,LoRA の低ランク射影と GPM 系の部分空間蓄積 \cite{gpm} によって同等の直交性を効率的に近似する点で実装上の差別化を図る.また本手法がリプレイ記憶を要したのに対し,提案手法は本体凍結+直交設計によりリハーサルなしでの ZCOCO 保持を狙う.

\subsection{比較:LoRA/PEFTベース継続学習手法の整理}

本章で扱った主要手法を,干渉抑制機構,パラメータ拡張の有無,タスク ID の要否,ゼロショット保持の実証,検出タスクへの対応,の観点で表~\ref{tab:lora_cl_compare} に整理する.

\begin{table}[t]
\centering
\caption{LoRA/PEFTベース継続学習手法の比較}
\label{tab:lora_cl_compare}
\begin{tabularx}{\linewidth}{l X c c c c}
\hline
手法 & 干渉抑制機構 & パラメータ & タスクID & ZS保持 & 検出対応 \\
\hline
LoRA \cite{lora} & なし(単体では非対応) & 一定 & 不要 & 未検証 & 容易 \\
AdaLoRA \cite{adalora} & なし(効率化が主目的) & 一定 & 不要 & 未検証 & 容易 \\
CODA-Prompt \cite{codaprompt} & プロンプト直交化+過去成分凍結 & 増(プール) & 暗黙に依存 & 未検証 & 限定的 \\
Low-rank Orth.\ \cite{lowrankorth} & Stiefel多様体上の直交部分空間 & 増(タスク毎) & 要(訓練時) & 未検証 & 限定的 \\
InfLoRA \cite{inflora} & $\mathcal{N}_t \cap \mathcal{M}_{t-1}^{\perp}$ への事前設計+DualGPM & 一定(統合) & 不要 & 未検証 & 適用可 \\
O-LoRA \cite{olora} & 直交損失 $\sum_i \|A_i^\top A_t\|_F^2$ & 増($A_i$保持) & 不要 & 実証(MMLU) & 適用可 \\
SD-LoRA \cite{sdlora} & 大きさ/方向分離+過去方向凍結 & 一定(非拡張) & 不要 & 未検証 & 適用可 \\
\hline
\end{tabularx}
\end{table}

表~\ref{tab:lora_cl_compare} の整理から,本研究の提案手法は,InfLoRA \cite{inflora} の干渉なし部分空間設計とパラメータ一定性を骨格とし,O-LoRA \cite{olora} が実証したゼロショット保持の発想を COCO(task0)への非対称な直交制約として取り込み,さらに SD-LoRA \cite{sdlora} のタスク ID 不要・非拡張性を継承する,という位置づけになる.既存手法はいずれも分類または言語タスクで検証されており,事前学習 VLM 検出器(MM-Grounding DINO)に対するドメイン増分かつ ZCOCO 保持という設定での適用は未開拓であり,ここに本研究の新規性がある.
